% Generated human-review packet template.  Source records, Lean statements,
% and raw-source-to-Spec screening fields are substituted by
% scripts/review_dashboard_packet.py.
% Do not hand-edit generated packets; annotate the PDF or reconcile notes into
% the paper-local human-review log.
\documentclass[10pt]{article}
\usepackage[margin=0.43in]{geometry}
\usepackage{fontspec}
\setmainfont{Latin Modern Roman}
\setmonofont{DejaVu Sans Mono}
\usepackage{hyperref}
\usepackage{amssymb}
\usepackage{fvextra}
\usepackage{enumitem}
\setlength{\parindent}{0pt}
\setlength{\parskip}{0.15em}
\linespread{0.92}
\hypersetup{colorlinks=true,linkcolor=blue,urlcolor=blue}
\DefineVerbatimEnvironment{ReviewVerbatim}{Verbatim}{breaklines=true,breakanywhere=true,fontsize=\scriptsize}
\newcommand{\reviewerbox}{%
  \par\noindent\fbox{\parbox{0.96\linewidth}{\rule{0pt}{0.38in}}}\par
}
\newcommand{\reviewmatch}[1]{%
  \par\noindent\textbf{Reviewer check:} $\square$\; #1\par
  \noindent\emph{If left unchecked, explain the discrepancy or uncertainty below.}\par
}

\begin{document}
\fontsize{9}{10}\selectfont

\begin{center}
{\LARGE Human review packet: DSWG24DiscretizationBias}\\[0.35em]
\small Generated from byte-pinned source excerpts, the expanded PaperInterface
specifications, and hash-bound raw-source-to-Spec screening records.
\end{center}

\noindent\textbf{Paper:} DSWG24DiscretizationBias\\
\textbf{Paper id:} \texttt{DSWG24DiscretizationBias}\\
\textbf{Source version:} arXiv:2405.16762; byte-pinned publication source record\\
\textbf{Source record:} \texttt{audit/paper\_statement\_map.json}\\
\textbf{Rows:} 17\\
\textbf{Generated:} 2026-08-19\\
\textbf{Source URL:} \href{https://arxiv.org/pdf/2405.16762}{official source}

\noindent\fbox{\parbox{0.96\linewidth}{\textbf{Review status:} 0/17 source-claim screens; 0/9 library prerequisite screens. This is a review worksheet while those direct checks remain pending; it is not a final validation receipt.}}\par



\section*{How to use this packet}
For each source claim, compare the exact source excerpts with the expanded
PaperInterface specification. Mark up this PDF or fill the blank reviewer box in the TeX source;
return the annotated file for reconciliation into the paper-local human-review
log.

\section*{Open the interactive dashboard}
The interactive dashboard is an optional alternative to using this PDF; it
presents the same review surface rather than an additional review requirement.
After forking and cloning (or pulling) EconCSLib, run the following from the
repository root. The cache command is needed only the first time in a new clone.
\begin{ReviewVerbatim}
cd <your-EconCSLib-checkout>
lake exe cache get
python3 scripts/review_dashboard.py --paper DSWG24DiscretizationBias --serve
\end{ReviewVerbatim}
Open \url{http://127.0.0.1:8765/} in a browser and keep that terminal running
while reviewing. The dashboard records saved annotations in this paper's
reviewer-owned review record.

\section*{Contents}
\begin{itemize}[leftmargin=1.2em]
\item \hyperlink{library-section-62111921f3668458}{Semantic library prerequisites (9; not paper claims)}
\begin{itemize}[leftmargin=1.2em]
\item \hyperlink{library-prerequisite-1ba21d6a145843fe}{EconCSLib.Decision.MonotoneExpectation}
\item \hyperlink{library-prerequisite-d23d58c11a4ba279}{EconCSLib.Decision.FiniteLinearExpectation}
\item \hyperlink{library-prerequisite-cf8d30fbaa7eab2a}{EconCSLib.Decision.IsPointwiseMax}
\item \hyperlink{library-prerequisite-aeeb723ddae150e9}{EconCSLib.Decision.averageScore}
\item \hyperlink{library-prerequisite-d8de4d5952ae79e0}{EconCSLib.Decision.datasetAccuracy}
\item \hyperlink{library-prerequisite-e0ab839a0c73984a}{EconCSLib.Decision.expectedDecisionAccuracy}
\item \hyperlink{library-prerequisite-5f9b15d72189cab5}{EconCSLib.Decision.expectedObjective}
\item \hyperlink{library-prerequisite-e5c7eedc40c11b7a}{EconCSLib.pmfExp}
\item \hyperlink{library-prerequisite-7fec9dedef722fc8}{EconCSLib.pmfProb}
\end{itemize}
\item \hyperlink{source-section-f10b5f68e4acf3dc}{Source claims (17)}
\begin{itemize}[leftmargin=1.2em]
\item \hyperlink{source-claim-1966e78967344232}{theorem1ii\_perfect\_classifier\_zero\_biasSpec}
\item \hyperlink{source-claim-94d82572ce30633e}{theorem2iii\_non\_argmax\_not\_paretoSpec}
\item \hyperlink{source-claim-9ffb824c48c4cf97}{theorem2iii\_weighted\_objective\_maximizer\_agrees\_argmaxSpec}
\item \hyperlink{source-claim-e99b78e6cf17d423}{tie\_broken\_argmax\_definitionSpec}
\item \hyperlink{source-claim-4bfc3aed59b0dc06}{thompson\_sampling\_definitionSpec}
\item \hyperlink{source-claim-4822dab29e865f72}{independent\_rule\_definitionSpec}
\item \hyperlink{source-claim-39ae869114aca78a}{integer\_optimization\_rule\_definitionSpec}
\item \hyperlink{source-claim-ec32b7cd29afa460}{bayes\_optimal\_definitionSpec}
\item \hyperlink{source-claim-7d95c2576dd31d3f}{pareto\_optimality\_definitionSpec}
\item \hyperlink{source-claim-4662dc8e8bb960d7}{nontrivial\_reference\_family\_definitionSpec}
\item \hyperlink{source-claim-d69492a6cc9555ec}{posterior\_simplex\_definitionSpec}
\item \hyperlink{source-claim-54a3ee9bbb8f086b}{calibration\_definitionSpec}
\item \hyperlink{source-claim-861b407dcf52f582}{theorem1i\_no\_information\_biasSpec}
\item \hyperlink{source-claim-d8ad786b7cc58022}{theorem1iii\_argmax\_bias\_le\_maeSpec}
\item \hyperlink{source-claim-d6b9878258465151}{theorem1iii\_tight\_binary\_exampleSpec}
\item \hyperlink{source-claim-82ddd45159ec7ca3}{theorem2i\_joint\_rule\_existsSpec}
\item \hyperlink{source-claim-cc8ba41c41752422}{theorem2ii\_argmax\_accuracy\_maximizingSpec}
\end{itemize}
\end{itemize}

\clearpage
\hypertarget{library-section-62111921f3668458}{}
\section*{Material library prerequisites}
\hypertarget{library-prerequisite-1ba21d6a145843fe}{}
\subsection*{\small\ttfamily\raggedright EconCSLib.\allowbreak{}Decision.\allowbreak{}MonotoneExpectation}
\textbf{Source connection:} no explicit byte-pinned paper source connection is registered.
\paragraph{Lean-expanded library semantic target}
\begin{ReviewVerbatim}
∀ {ω : Type u_1} (expect : (ω → ℝ) → ℝ) (f g : ω → ℝ), (∀ (x : ω), f x ≤ g x) → expect f ≤ expect g
\end{ReviewVerbatim}

\paragraph{Exact Lean library declaration}
\begin{ReviewVerbatim}
def MonotoneExpectation {ω : Type*} (expect : (ω → ℝ) → ℝ) : Prop :=
  ∀ f g, (∀ x, f x ≤ g x) → expect f ≤ expect g
\end{ReviewVerbatim}

\paragraph{Recorded source-to-library screening}
\textbf{Verdict:} \texttt{not recorded}
\reviewmatch{Matches source input}
\noindent\textbf{Reviewer annotation}\par
\reviewerbox
\clearpage
\hypertarget{library-prerequisite-d23d58c11a4ba279}{}
\subsection*{\small\ttfamily\raggedright EconCSLib.\allowbreak{}Decision.\allowbreak{}FiniteLinearExpectation}
\textbf{Source connection:} no explicit byte-pinned paper source connection is registered.
\paragraph{Lean-expanded library semantic target}
\begin{ReviewVerbatim}
∀ {ω : Type u_1} (expect : (ω → ℝ) → ℝ),
  EconCSLib.Decision.MonotoneExpectation expect ∧
    (∀ (c : ℝ) (f : ω → ℝ), (expect fun x => c * f x) = c * expect f) ∧
      ∀ (f g : ω → ℝ), (expect fun x => f x + g x) = expect f + expect g
\end{ReviewVerbatim}

\paragraph{Exact Lean library declaration}
\begin{ReviewVerbatim}
def FiniteLinearExpectation {ω : Type*} (expect : (ω → ℝ) → ℝ) : Prop :=
  MonotoneExpectation expect ∧
    (∀ c f, expect (fun x => c * f x) = c * expect f) ∧
    (∀ f g, expect (fun x => f x + g x) = expect f + expect g)
\end{ReviewVerbatim}

\paragraph{Recorded source-to-library screening}
\textbf{Verdict:} \texttt{not recorded}
\reviewmatch{Matches source input}
\noindent\textbf{Reviewer annotation}\par
\reviewerbox
\clearpage
\hypertarget{library-prerequisite-cf8d30fbaa7eab2a}{}
\subsection*{\small\ttfamily\raggedright EconCSLib.\allowbreak{}Decision.\allowbreak{}IsPointwiseMax}
\textbf{Source connection:} no explicit byte-pinned paper source connection is registered.
\paragraph{Lean-expanded library semantic target}
\begin{ReviewVerbatim}
∀ {ι : Type u_1} {α : Type u_2} (score : ι → α → ℝ) (choose : ι → α) (i : ι) (a : α), score i a ≤ score i (choose i)
\end{ReviewVerbatim}

\paragraph{Exact Lean library declaration}
\begin{ReviewVerbatim}
def IsPointwiseMax {ι α : Type*} (score : ι → α → ℝ) (choose : ι → α) : Prop :=
  ∀ i a, score i a ≤ score i (choose i)
\end{ReviewVerbatim}

\paragraph{Recorded source-to-library screening}
\textbf{Verdict:} \texttt{not recorded}
\reviewmatch{Matches source input}
\noindent\textbf{Reviewer annotation}\par
\reviewerbox
\clearpage
\hypertarget{library-prerequisite-aeeb723ddae150e9}{}
\subsection*{\small\ttfamily\raggedright EconCSLib.\allowbreak{}Decision.\allowbreak{}averageScore}
\textbf{Source connection:} no explicit byte-pinned paper source connection is registered.
\paragraph{Lean-expanded library semantic target}
\begin{ReviewVerbatim}
{ι : Type u_1} →
  {α : Type u_2} →
    [inst : Fintype ι] → (score : ι → α → ℝ) → (choose : ι → α) → (∑ i, score i (choose i)) / ↑(Fintype.card ι)
\end{ReviewVerbatim}

\paragraph{Exact Lean library declaration}
\begin{ReviewVerbatim}
noncomputable def averageScore {ι α : Type*} [Fintype ι]
    (score : ι → α → ℝ) (choose : ι → α) : ℝ :=
  (∑ i : ι, score i (choose i)) / (Fintype.card ι : ℝ)
\end{ReviewVerbatim}

\paragraph{Recorded source-to-library screening}
\textbf{Verdict:} \texttt{not recorded}
\reviewmatch{Matches source input}
\noindent\textbf{Reviewer annotation}\par
\reviewerbox
\clearpage
\hypertarget{library-prerequisite-d8de4d5952ae79e0}{}
\subsection*{\small\ttfamily\raggedright EconCSLib.\allowbreak{}Decision.\allowbreak{}datasetAccuracy}
\textbf{Source connection:} no explicit byte-pinned paper source connection is registered.
\paragraph{Lean-expanded library semantic target}
\begin{ReviewVerbatim}
{ι : Type u_1} →
  {α : Type u_2} →
    [inst : Fintype ι] →
      [inst_1 : DecidableEq α] →
        (truth decision : ι → α) → EconCSLib.Decision.averageScore (fun i a => if a = truth i then 1 else 0) decision
\end{ReviewVerbatim}

\paragraph{Exact Lean library declaration}
\begin{ReviewVerbatim}
noncomputable def datasetAccuracy {ι α : Type*} [Fintype ι] [DecidableEq α]
    (truth decision : ι → α) : ℝ :=
  averageScore (fun i a => if a = truth i then (1 : ℝ) else 0) decision
\end{ReviewVerbatim}

\paragraph{Recorded source-to-library screening}
\textbf{Verdict:} \texttt{not recorded}
\reviewmatch{Matches source input}
\noindent\textbf{Reviewer annotation}\par
\reviewerbox
\clearpage
\hypertarget{library-prerequisite-e0ab839a0c73984a}{}
\subsection*{\small\ttfamily\raggedright EconCSLib.\allowbreak{}Decision.\allowbreak{}expectedDecisionAccuracy}
\textbf{Source connection:} no explicit byte-pinned paper source connection is registered.
\paragraph{Lean-expanded library semantic target}
\begin{ReviewVerbatim}
{ω : Type u_1} →
  {σ : Type u_2} →
    {ι : Type u_3} →
      {α : Type u_4} →
        [inst : Fintype ι] →
          [inst_1 : DecidableEq α] →
            (expect : (ω → ℝ) → ℝ) →
              (obs : ω → σ) →
                (truth : ω → ι → α) →
                  (rule : σ → ι → α) → expect fun x => EconCSLib.Decision.datasetAccuracy (truth x) (rule (obs x))
\end{ReviewVerbatim}

\paragraph{Exact Lean library declaration}
\begin{ReviewVerbatim}
noncomputable def expectedDecisionAccuracy {ω σ ι α : Type*}
    [Fintype ι] [DecidableEq α]
    (expect : (ω → ℝ) → ℝ) (obs : ω → σ)
    (truth : ω → ι → α) (rule : σ → ι → α) : ℝ :=
  expect (fun x => datasetAccuracy (truth x) (rule (obs x)))
\end{ReviewVerbatim}

\paragraph{Recorded source-to-library screening}
\textbf{Verdict:} \texttt{not recorded}
\reviewmatch{Matches source input}
\noindent\textbf{Reviewer annotation}\par
\reviewerbox
\clearpage
\hypertarget{library-prerequisite-5f9b15d72189cab5}{}
\subsection*{\small\ttfamily\raggedright EconCSLib.\allowbreak{}Decision.\allowbreak{}expectedObjective}
\textbf{Source connection:} no explicit byte-pinned paper source connection is registered.
\paragraph{Lean-expanded library semantic target}
\begin{ReviewVerbatim}
{ω : Type u_1} →
  {σ : Type u_2} →
    {β : Type u_3} →
      (expect : (ω → ℝ) → ℝ) →
        (obs : ω → σ) → (objective : σ → β → ℝ) → (rule : σ → β) → expect fun x => objective (obs x) (rule (obs x))
\end{ReviewVerbatim}

\paragraph{Exact Lean library declaration}
\begin{ReviewVerbatim}
noncomputable def expectedObjective {ω σ β : Type*}
    (expect : (ω → ℝ) → ℝ) (obs : ω → σ)
    (objective : σ → β → ℝ) (rule : σ → β) : ℝ :=
  expect (fun x => objective (obs x) (rule (obs x)))
\end{ReviewVerbatim}

\paragraph{Recorded source-to-library screening}
\textbf{Verdict:} \texttt{not recorded}
\reviewmatch{Matches source input}
\noindent\textbf{Reviewer annotation}\par
\reviewerbox
\clearpage
\hypertarget{library-prerequisite-e5c7eedc40c11b7a}{}
\subsection*{\small\ttfamily\raggedright EconCSLib.\allowbreak{}pmfExp}
\textbf{Source connection:} no explicit byte-pinned paper source connection is registered.
\paragraph{Lean-expanded library semantic target}
\begin{ReviewVerbatim}
{α : Type u_1} → [inst : Fintype α] → [DecidableEq α] → (μ : PMF α) → (f : α → ℝ) → ∑ a, (μ a).toReal * f a
\end{ReviewVerbatim}

\paragraph{Exact Lean library declaration}
\begin{ReviewVerbatim}
noncomputable def pmfExp {α : Type*} [Fintype α] [DecidableEq α]
    (μ : PMF α) (f : α → ℝ) : ℝ :=
  ∑ a, (μ a).toReal * f a
\end{ReviewVerbatim}

\paragraph{Recorded source-to-library screening}
\textbf{Verdict:} \texttt{not recorded}
\reviewmatch{Matches source input}
\noindent\textbf{Reviewer annotation}\par
\reviewerbox
\clearpage
\hypertarget{library-prerequisite-7fec9dedef722fc8}{}
\subsection*{\small\ttfamily\raggedright EconCSLib.\allowbreak{}pmfProb}
\textbf{Source connection:} no explicit byte-pinned paper source connection is registered.
\paragraph{Lean-expanded library semantic target}
\begin{ReviewVerbatim}
{α : Type u_1} →
  [inst : Fintype α] →
    [inst_1 : DecidableEq α] →
      (μ : PMF α) → (p : α → Prop) → [inst_2 : DecidablePred p] → EconCSLib.pmfExp μ fun a => if p a then 1 else 0
\end{ReviewVerbatim}

\paragraph{Exact Lean library declaration}
\begin{ReviewVerbatim}
noncomputable def pmfProb {α : Type*} [Fintype α] [DecidableEq α]
    (μ : PMF α) (p : α → Prop) [DecidablePred p] : ℝ :=
  pmfExp μ (fun a => if p a then 1 else 0)
\end{ReviewVerbatim}

\paragraph{Recorded source-to-library screening}
\textbf{Verdict:} \texttt{not recorded}
\reviewmatch{Matches source input}
\noindent\textbf{Reviewer annotation}\par
\reviewerbox

\clearpage
\hypertarget{source-claim-1966e78967344232}{}
\hypertarget{source-section-f10b5f68e4acf3dc}{}
\section*{Source claims}
\subsection*{Review row 1}
\textbf{Source locator:} {\footnotesize\raggedright cited publication, lines 642-\allowbreak{}666\par}
\paragraph{Verbatim source input}
\begin{ReviewVerbatim}
Theorem 1. Argmax bias depends on predictive uncertainty, i.e., how much information features x provide
about the true class label. Consider calibrated classifier q and the argmax decision rule Dargmax . Let N > K
and consider a reference distribution of either the aggregate posterior or the prior. Then,
  (i). Suppose x provides no information and q(y, x) = Pr(y) for all x, y. Then, amplification bias is
       maximized and all data points are classified as a plurality class:
                                                  (
                                                    1 − Pr(y), if y = arg maxw p(w),
                            BIAS(y, Dargmax ) =
                                                    − Pr(y),      otherwise.
   9 Note that, though related, mean average error of the continuous predictor q(y, x) is not equivalent to expected zero-one

loss of the discrete decisions. The latter, our measure of decision accuracy, also depends on how the continuous predictor is
discretized.                                n                                           o
  10 For all c of non-zero measure: ∀c ∈
                                                 R
                                             c : (x,y) I[q(y, x) = c]f (x, y)d(y, x) > 0 . Further note that calibration implies
Bayesian consistency, i.e., the continuous model cumulatively assigns the proper mass to each class y: EX [q(y, x)] = Pr(y).


                                                              11

[form-feed in source extraction]
                                                                                   (
                                                                                    1,                                                                        if z = y,
  (ii). Suppose the classifier is perfect, i.e., ∀(x, y) ∼ FXY , we have q(z, x) =                                                                                       Then, there
                                                                                    0,                                                                        otherwise.
        is no amplification bias, BIAS(y, Dargmax ) = 0.

[Next verbatim source excerpt]

Theorem 1. Argmax bias depends on predictive uncertainty, i.e., how much information features x provide
about the true class label. Consider calibrated classifier q and the argmax decision rule Dargmax . Let N > K
and consider a reference distribution of either the aggregate posterior or the prior. Then,
  (i). Suppose x provides no information and q(y, x) = Pr(y) for all x, y. Then, amplification bias is
       maximized and all data points are classified as a plurality class:
                                                  (
                                                    1 − Pr(y), if y = arg maxw p(w),
                            BIAS(y, Dargmax ) =
                                                    − Pr(y),      otherwise.
   9 Note that, though related, mean average error of the continuous predictor q(y, x) is not equivalent to expected zero-one

loss of the discrete decisions. The latter, our measure of decision accuracy, also depends on how the continuous predictor is
discretized.                                n                                           o
  10 For all c of non-zero measure: ∀c ∈
                                                 R
                                             c : (x,y) I[q(y, x) = c]f (x, y)d(y, x) > 0 . Further note that calibration implies
Bayesian consistency, i.e., the continuous model cumulatively assigns the proper mass to each class y: EX [q(y, x)] = Pr(y).


                                                              11

[form-feed in source extraction]
                                                                                   (
                                                                                    1,                                                                        if z = y,
  (ii). Suppose the classifier is perfect, i.e., ∀(x, y) ∼ FXY , we have q(z, x) =                                                                                       Then, there
                                                                                    0,                                                                        otherwise.
        is no amplification bias, BIAS(y, Dargmax ) = 0.
\end{ReviewVerbatim}


\paragraph{Expanded PaperInterface specification}
\begin{ReviewVerbatim}
∀ {X : Type u_1} {Y : Type u_2} [inst : Fintype X] [inst_1 : DecidableEq X] [inst_2 : Fintype Y]
  [inst_3 : DecidableEq Y] (μ : PMF (X × Y)) (q : X → Y → ℝ) (rule : X → Y) (y : Y),
  DSWG24DiscretizationBias.ProofBridge.perfectClassifierOnSupport μ q rule →
    DSWG24DiscretizationBias.Finite.paperBias μ rule (DSWG24DiscretizationBias.ProofBridge.prior μ) y = 0 ∧
      DSWG24DiscretizationBias.Finite.paperBias μ rule (DSWG24DiscretizationBias.Finite.paperAggregatePosterior μ q) y =
        0
\end{ReviewVerbatim}

\paragraph{Lean proof endpoint (built separately)}\begin{ReviewVerbatim}
DSWG24DiscretizationBias.PaperInterface.theorem1ii_perfect_classifier_zero_bias
\end{ReviewVerbatim}

\paragraph{Recorded source-to-Spec screening}
\textbf{Verdict:} \texttt{not recorded}
\\
\textbf{Reason:} not recorded
\reviewmatch{Matches source input}
\noindent\textbf{Reviewer annotation}\par
\reviewerbox
\clearpage
\hypertarget{source-claim-94d82572ce30633e}{}
\section*{Review row 2}
\textbf{Source locator:} {\footnotesize\raggedright cited publication, lines 756-\allowbreak{}763\par}
\paragraph{Verbatim source input}
\begin{ReviewVerbatim}
Theorem 2. Consider Bayes optimal q, and N > K. For all F :
   (i). For every γ, pref there exists a joint decision-making rule that maximizes Oγ (D, q, pref ) in Equa-
        tion (3).

  (ii). The argmax decision rule Dargmax maximizes O1 (D, q, pref ), i.e., is accuracy maximizing.
 (iii). Dargmax is the only Pareto optimal independent decision rule for non-trivial pref ∈ P. No independent
        decision rule D maximizes objective Oγ for any γ, unless γ = 1 and D agrees with Dargmax with
        probability 1.
\end{ReviewVerbatim}


\paragraph{Expanded PaperInterface specification}
\begin{ReviewVerbatim}
∀ {X : Type u_1} [inst : Fintype X] [inst_1 : DecidableEq X] {N K : ℕ} (μ : PMF X) (posterior : X → Fin K → ℝ)
  (rule argmaxRule : X → Fin K) (aggregateArgmax : (Fin N → X) → Fin K) (prefAt : (Fin N → X) → Fin K → ℝ),
  (∀ (x : X), DSWG24DiscretizationBias.ProofBridge.isFirstArgmax (posterior x) (argmaxRule x)) →
    DSWG24DiscretizationBias.ProofBridge.sourcePNq
        (fun sample => DSWG24DiscretizationBias.Theorem2iii.sampledPosterior posterior sample) aggregateArgmax prefAt →
      0 < N →
        (DSWG24DiscretizationBias.Pareto.ParetoOptimal
            (DSWG24DiscretizationBias.Theorem2iii.expectedDatasetAccuracy
              (DSWG24DiscretizationBias.Theorem2iii.iidSamplePMF μ N) fun sample =>
              DSWG24DiscretizationBias.Theorem2iii.sampledPosterior posterior sample)
            (DSWG24DiscretizationBias.Theorem2iii.expectedDatasetFidelityAt
              (DSWG24DiscretizationBias.Theorem2iii.iidSamplePMF μ N) prefAt)
            fun sample => DSWG24DiscretizationBias.Theorem2iii.sampledDecision rule sample) →
          (EconCSLib.pmfProb μ fun x => rule x ≠ argmaxRule x) = 0
\end{ReviewVerbatim}

\paragraph{Lean proof endpoint (built separately)}\begin{ReviewVerbatim}
DSWG24DiscretizationBias.PaperInterface.theorem2iii_non_argmax_not_pareto
\end{ReviewVerbatim}

\paragraph{Recorded source-to-Spec screening}
\textbf{Verdict:} \texttt{not recorded}
\\
\textbf{Reason:} not recorded
\reviewmatch{Matches source input}
\noindent\textbf{Reviewer annotation}\par
\reviewerbox
\clearpage
\hypertarget{source-claim-9ffb824c48c4cf97}{}
\section*{Review row 3}
\textbf{Source locator:} {\footnotesize\raggedright cited publication, lines 756-\allowbreak{}763\par}
\paragraph{Verbatim source input}
\begin{ReviewVerbatim}
Theorem 2. Consider Bayes optimal q, and N > K. For all F :
   (i). For every γ, pref there exists a joint decision-making rule that maximizes Oγ (D, q, pref ) in Equa-
        tion (3).

  (ii). The argmax decision rule Dargmax maximizes O1 (D, q, pref ), i.e., is accuracy maximizing.
 (iii). Dargmax is the only Pareto optimal independent decision rule for non-trivial pref ∈ P. No independent
        decision rule D maximizes objective Oγ for any γ, unless γ = 1 and D agrees with Dargmax with
        probability 1.
\end{ReviewVerbatim}


\paragraph{Expanded PaperInterface specification}
\begin{ReviewVerbatim}
∀ {X : Type u_1} [inst : Fintype X] [inst_1 : DecidableEq X] {N K : ℕ} (μ : PMF X) (posterior : X → Fin K → ℝ)
  (rule argmaxRule : X → Fin K) (aggregateArgmax : (Fin N → X) → Fin K) (prefAt : (Fin N → X) → Fin K → ℝ) (γ : ℝ),
  (∀ (x : X), DSWG24DiscretizationBias.ProofBridge.isFirstArgmax (posterior x) (argmaxRule x)) →
    DSWG24DiscretizationBias.ProofBridge.sourcePNq
        (fun sample => DSWG24DiscretizationBias.Theorem2iii.sampledPosterior posterior sample) aggregateArgmax prefAt →
      0 ≤ γ →
        γ < 1 →
          0 < N →
            (∀ (other : (Fin N → X) → Fin N → Fin K),
                DSWG24DiscretizationBias.Pareto.weightedObjective γ
                    (DSWG24DiscretizationBias.Theorem2iii.expectedDatasetAccuracy
                      (DSWG24DiscretizationBias.Theorem2iii.iidSamplePMF μ N) fun sample =>
                      DSWG24DiscretizationBias.Theorem2iii.sampledPosterior posterior sample)
                    (DSWG24DiscretizationBias.Theorem2iii.expectedDatasetFidelityAt
                      (DSWG24DiscretizationBias.Theorem2iii.iidSamplePMF μ N) prefAt)
                    other ≤
                  DSWG24DiscretizationBias.Pareto.weightedObjective γ
                    (DSWG24DiscretizationBias.Theorem2iii.expectedDatasetAccuracy
                      (DSWG24DiscretizationBias.Theorem2iii.iidSamplePMF μ N) fun sample =>
                      DSWG24DiscretizationBias.Theorem2iii.sampledPosterior posterior sample)
                    (DSWG24DiscretizationBias.Theorem2iii.expectedDatasetFidelityAt
                      (DSWG24DiscretizationBias.Theorem2iii.iidSamplePMF μ N) prefAt)
                    fun sample => DSWG24DiscretizationBias.Theorem2iii.sampledDecision rule sample) →
              (EconCSLib.pmfProb μ fun x => rule x ≠ argmaxRule x) = 0
\end{ReviewVerbatim}

\paragraph{Lean proof endpoint (built separately)}\begin{ReviewVerbatim}
DSWG24DiscretizationBias.PaperInterface.theorem2iii_weighted_objective_maximizer_agrees_argmax
\end{ReviewVerbatim}

\paragraph{Recorded source-to-Spec screening}
\textbf{Verdict:} \texttt{not recorded}
\\
\textbf{Reason:} not recorded
\reviewmatch{Matches source input}
\noindent\textbf{Reviewer annotation}\par
\reviewerbox
\clearpage
\hypertarget{source-claim-e99b78e6cf17d423}{}
\section*{Review row 4}
\textbf{Source locator:} {\footnotesize\raggedright cited publication, lines 215-\allowbreak{}231\par}
\paragraph{Verbatim source input}
\begin{ReviewVerbatim}
1.1.2    Overview of decision rules
Several rules are used in academia and in practice:

Argmax rule assigns, for each data point, the most likely2 class: Dargmax ({q(xi }) ≜ {ŷi = arg maxy q(y, xi )}.
Threshold at t rule assigns a label only if the argmax class has a probability of at least t, i.e., ŷi =
    arg maxy q(y, xi ) if maxy q(y, xi ) ≥ t, otherwise i is left uncoded.

Thompson sampling for each data point i, samples a class y based on the probabilities q(y, xi ), i.e., ŷi = y
    with probability q(y, xi ), and so in expectation will have labels matching the aggregate posterior.

    These rules are all independent: the assigned label for a data point i depends on q(xi ) but not on other
data points or their probability vectors {q(xj )}j̸=i .
    In contrast, our proposed approach also includes joint decision optimization rules, which depend on the
entire set of data points, as a solution to discretization bias.
   2 Suppose ties are broken according to some consistent ordering over the classes.      In the event of ties, we will use
arg maxy q(y, x) to denote the single class that is in the argmax set and is first in the consistent ordering among classes
in the argmax set.

[Next verbatim source excerpt]

Argmax rule assigns, for each data point, the most likely2 class: Dargmax ({q(xi }) ≜ {ŷi = arg maxy q(y, xi )}.
Threshold at t rule assigns a label only if the argmax class has a probability of at least t, i.e., ŷi =
    arg maxy q(y, xi ) if maxy q(y, xi ) ≥ t, otherwise i is left uncoded.

Thompson sampling for each data point i, samples a class y based on the probabilities q(y, xi ), i.e., ŷi = y
    with probability q(y, xi ), and so in expectation will have labels matching the aggregate posterior.

    These rules are all independent: the assigned label for a data point i depends on q(xi ) but not on other
data points or their probability vectors {q(xj )}j̸=i .
    In contrast, our proposed approach also includes joint decision optimization rules, which depend on the
entire set of data points, as a solution to discretization bias.
   2 Suppose ties are broken according to some consistent ordering over the classes.      In the event of ties, we will use
arg maxy q(y, x) to denote the single class that is in the argmax set and is first in the consistent ordering among classes
in the argmax set.
\end{ReviewVerbatim}


\paragraph{Expanded PaperInterface specification}
\begin{ReviewVerbatim}
∀ {N K : ℕ} (q : Fin N → Fin K → ℝ) (rule : Fin N → Fin K),
  DSWG24DiscretizationBias.ProofBridge.isTieBrokenArgmaxRule q rule ↔
    ∀ (i : Fin N), DSWG24DiscretizationBias.ProofBridge.isFirstArgmax (q i) (rule i)
\end{ReviewVerbatim}

\paragraph{Lean proof endpoint (built separately)}\begin{ReviewVerbatim}
DSWG24DiscretizationBias.PaperInterface.tie_broken_argmax_definition
\end{ReviewVerbatim}

\paragraph{Recorded source-to-Spec screening}
\textbf{Verdict:} \texttt{not recorded}
\\
\textbf{Reason:} not recorded
\reviewmatch{Matches source input}
\noindent\textbf{Reviewer annotation}\par
\reviewerbox
\clearpage
\hypertarget{source-claim-4bfc3aed59b0dc06}{}
\section*{Review row 5}
\textbf{Source locator:} {\footnotesize\raggedright cited publication, lines 222-\allowbreak{}223\par}
\paragraph{Verbatim source input}
\begin{ReviewVerbatim}
Thompson sampling for each data point i, samples a class y based on the probabilities q(y, xi ), i.e., ŷi = y
    with probability q(y, xi ), and so in expectation will have labels matching the aggregate posterior.
\end{ReviewVerbatim}


\paragraph{Expanded PaperInterface specification}
\begin{ReviewVerbatim}
∀ {N K : ℕ} (q selectionProbability : Fin N → Fin K → ℝ),
  DSWG24DiscretizationBias.ProofBridge.isThompsonSamplingRule q selectionProbability ↔
    ∀ (i : Fin N) (y : Fin K), selectionProbability i y = q i y
\end{ReviewVerbatim}

\paragraph{Lean proof endpoint (built separately)}\begin{ReviewVerbatim}
DSWG24DiscretizationBias.PaperInterface.thompson_sampling_definition
\end{ReviewVerbatim}

\paragraph{Recorded source-to-Spec screening}
\textbf{Verdict:} \texttt{not recorded}
\\
\textbf{Reason:} not recorded
\reviewmatch{Matches source input}
\noindent\textbf{Reviewer annotation}\par
\reviewerbox
\clearpage
\hypertarget{source-claim-4822dab29e865f72}{}
\section*{Review row 6}
\textbf{Source locator:} {\footnotesize\raggedright cited publication, lines 225-\allowbreak{}228\par}
\paragraph{Verbatim source input}
\begin{ReviewVerbatim}
These rules are all independent: the assigned label for a data point i depends on q(xi ) but not on other
data points or their probability vectors {q(xj )}j̸=i .
    In contrast, our proposed approach also includes joint decision optimization rules, which depend on the
entire set of data points, as a solution to discretization bias.

[Next verbatim source excerpt]

data point xi . A (potentially randomized) decision (discretization) rule D : △NY →Y
                                                                                      N
                                                                                         is a rule that, given a
                                   N
set of probability vectors {q(xi )}i=1 associated with N data points, assigns each data point i a label ŷi ∈ Y.

[Next verbatim source excerpt]

    Decision (discretization) rules are as defined in Section 1.1. Our theoretical results distinguish between
independent (e.g., argmax, threhsolding, Thompson sampling) and joint (e.g., our optimization approach)
rules. Formally, with independent rules, there exists a (potentially randomized) function d such that
Dd ({q(xi )}) = {ŷi = d(q(y, xi ))}.
\end{ReviewVerbatim}


\paragraph{Expanded PaperInterface specification}
\begin{ReviewVerbatim}
∀ {X : Type u_1} {N K : ℕ} (rule : (Fin N → X) → Fin N → Fin K),
  DSWG24DiscretizationBias.ProofBridge.isIndependentRule rule ↔
    ∃ d, ∀ (xs : Fin N → X) (i : Fin N), rule xs i = d (xs i)
\end{ReviewVerbatim}

\paragraph{Lean proof endpoint (built separately)}\begin{ReviewVerbatim}
DSWG24DiscretizationBias.PaperInterface.independent_rule_definition
\end{ReviewVerbatim}

\paragraph{Recorded source-to-Spec screening}
\textbf{Verdict:} \texttt{not recorded}
\\
\textbf{Reason:} not recorded
\reviewmatch{Matches source input}
\noindent\textbf{Reviewer annotation}\par
\reviewerbox
\clearpage
\hypertarget{source-claim-39ae869114aca78a}{}
\section*{Review row 7}
\textbf{Source locator:} {\footnotesize\raggedright cited publication, lines 252-\allowbreak{}264\par}
\paragraph{Verbatim source input}
\begin{ReviewVerbatim}
Integer optimization rules directly optimize given objectives. We consider the rule that corresponds to
     solving the following optimization problem, for a given γ, data points {xi }, and reference pref :
                                                            N
                                                                           !
                      γ                                  1 X
                   D ({q(xi )}, pref ) = arg max{ŷi } γ       q(ŷi , xi ) + (1 − γ)fid(pref , ŷ1:N ). (1)
                                                         N i=1

     The rule balances, parameterized by γ ∈ [0, 1], row-level scores (accuracy as measured by q) and
     distributional fidelity.
Aggregate posterior matching A special case of the joint optimization decision rules defined in Equa-
    tion (1) is as γ → 0, leading to a rule that maximizes accuracy subject to the constraint that distri-
    butional fidelity be as high as possible. We call such rules matching solutions, e.g., aggregate posterior
\end{ReviewVerbatim}


\paragraph{Expanded PaperInterface specification}
\begin{ReviewVerbatim}
∀ {N K : ℕ} (γ : ℝ) (q : Fin N → Fin K → ℝ) (pref : Fin K → ℝ) (decision : Fin N → Fin K),
  DSWG24DiscretizationBias.ProofBridge.maximizesEquation1 γ q pref decision ↔
    ∀ (other : Fin N → Fin K),
      DSWG24DiscretizationBias.ProofBridge.equation1Objective γ q pref other ≤
        DSWG24DiscretizationBias.ProofBridge.equation1Objective γ q pref decision
\end{ReviewVerbatim}

\paragraph{Lean proof endpoint (built separately)}\begin{ReviewVerbatim}
DSWG24DiscretizationBias.PaperInterface.integer_optimization_rule_definition
\end{ReviewVerbatim}

\paragraph{Recorded source-to-Spec screening}
\textbf{Verdict:} \texttt{not recorded}
\\
\textbf{Reason:} not recorded
\reviewmatch{Matches source input}
\noindent\textbf{Reviewer annotation}\par
\reviewerbox
\clearpage
\hypertarget{source-claim-ec32b7cd29afa460}{}
\section*{Review row 8}
\textbf{Source locator:} {\footnotesize\raggedright cited publication, lines 597-\allowbreak{}607\par}
\paragraph{Verbatim source input}
\begin{ReviewVerbatim}
3.1.1    Continuous classifier and discretization rules
Suppose we have a machine learning classifier q (e.g., trained on some other dataset) that makes continuous
probability predictions for each class, given the data. Let q(y, x) denote the probability placed on class y



                                                       10

[form-feed in source extraction]
given data x. For example, the Bayes optimal classifier is
                                               q ∗ (y, x) ≜ Pr(Y = y|X = x).
We will refer to p as the prior and q as our learned posterior. Let the simplex of possible probability vectors
over the support Y be △Y , and let q(x) ∈ △Y be the vector of probabilities (q(y, x))y .
\end{ReviewVerbatim}


\paragraph{Expanded PaperInterface specification}
\begin{ReviewVerbatim}
∀ {Ω : Type u_1} {σ : Type u_2} [inst : Fintype Ω] [inst_1 : DecidableEq Ω] [inst_2 : Fintype σ]
  [inst_3 : DecidableEq σ] {N K : ℕ} (μ : PMF Ω) (observedDataset : Ω → σ) (trueLabels : Ω → Fin N → Fin K)
  (posterior : σ → Fin N → Fin K → ℝ),
  DSWG24DiscretizationBias.ProofBridge.bayesOptimal μ observedDataset trueLabels posterior ↔
    ∀ (xs : σ) (i : Fin N) (y : Fin K),
      (EconCSLib.pmfProb μ fun w => observedDataset w = xs ∧ trueLabels w i = y) =
        posterior xs i y * EconCSLib.pmfProb μ fun w => observedDataset w = xs
\end{ReviewVerbatim}

\paragraph{Lean proof endpoint (built separately)}\begin{ReviewVerbatim}
DSWG24DiscretizationBias.PaperInterface.bayes_optimal_definition
\end{ReviewVerbatim}

\paragraph{Recorded source-to-Spec screening}
\textbf{Verdict:} \texttt{not recorded}
\\
\textbf{Reason:} not recorded
\reviewmatch{Matches source input}
\noindent\textbf{Reviewer annotation}\par
\reviewerbox
\clearpage
\hypertarget{source-claim-7d95c2576dd31d3f}{}
\section*{Review row 9}
\textbf{Source locator:} {\footnotesize\raggedright cited publication, lines 620-\allowbreak{}633\par}
\paragraph{Verbatim source input}
\begin{ReviewVerbatim}
3.1.2    Formal metrics and Pareto optimal decision-making
For our theoretical results, we consider the expectations of our accuracy, fidelity, and bias metrics, calculated
over the joint distribution of the data FXY over datasets of size N (and any randomness in the decision
rule). The expected accuracy of a decision rule D together with continuous predictor q is thus
                                     ACCN (D, q) = EFXY [acc(y1:N , D({q(xi )}))] .                                            (2)
Analogously FIDN (D, q, pref ) is the expectation of fidelity and BIASN (y, D, q, pref ) is the expectation of bias.
   For a given data distribution F , continuous classifier q, reference distribution pref , and dataset size, we
want to make effective decisions. Given the multiple desiderata, we thus want Pareto optimal decision-
making, i.e., to adopt a rule D that maximizes, for some γ ∈ (0, 1],
                              γ
                             ON (D, q, pref ) = γACCN (D, q) + (1 − γ)FIDN (D, q, pref ).                                      (3)
The weight γ is task-dependent, set by a practitioner’s expertise for whether accuracy or distributional
fidelity is more important.
    When clear from context, we omit q, pref , and N from the arguments for BIAS, Oγ , ACC, and FID.
\end{ReviewVerbatim}


\paragraph{Expanded PaperInterface specification}
\begin{ReviewVerbatim}
∀ {Rule : Type u_1} (accuracyMetric fidelityMetric : Rule → ℝ) (rule : Rule),
  DSWG24DiscretizationBias.ProofBridge.paretoOptimal accuracyMetric fidelityMetric rule ↔
    DSWG24DiscretizationBias.Pareto.ParetoOptimal accuracyMetric fidelityMetric rule
\end{ReviewVerbatim}

\paragraph{Lean proof endpoint (built separately)}\begin{ReviewVerbatim}
DSWG24DiscretizationBias.PaperInterface.pareto_optimality_definition
\end{ReviewVerbatim}

\paragraph{Recorded source-to-Spec screening}
\textbf{Verdict:} \texttt{not recorded}
\\
\textbf{Reason:} not recorded
\reviewmatch{Matches source input}
\noindent\textbf{Reviewer annotation}\par
\reviewerbox
\clearpage
\hypertarget{source-claim-4662dc8e8bb960d7}{}
\section*{Review row 10}
\textbf{Source locator:} {\footnotesize\raggedright cited publication, lines 744-\allowbreak{}755\par}
\paragraph{Verbatim source input}
\begin{ReviewVerbatim}
3.3     Results: Individual discretization versus joint optimization
Above, we study how the continuous classifier’s properties induce discretization bias, when using the argmax
rule. We now study discrete decision-making: given an imperfect but Bayes optimal classifier, how can
we make unbiased decisions? We show that joint decision-making, as in the program in Equation (1), is
necessary for Pareto optimality.
                                                                                     q
    For the result, we define the notion of non-trivial reference distributions PN     : a reference distribution
                                                       1
pq : {xi } → △NY is in P N if it assigns mass at least N to the most likely class in the  dataset, i.e., at least 1
data point would be discretized to the most likely class arg maxy i q(y, xi ).13
                                                                   P


[Next verbatim source excerpt]

3.3     Results: Individual discretization versus joint optimization
Above, we study how the continuous classifier’s properties induce discretization bias, when using the argmax
rule. We now study discrete decision-making: given an imperfect but Bayes optimal classifier, how can
we make unbiased decisions? We show that joint decision-making, as in the program in Equation (1), is
necessary for Pareto optimality.
                                                                                     q
    For the result, we define the notion of non-trivial reference distributions PN     : a reference distribution
                                                       1
pq : {xi } → △NY is in P N if it assigns mass at least N to the most likely class in the  dataset, i.e., at least 1
data point would be discretized to the most likely class arg maxy i q(y, xi ).13
                                                                   P


[Next verbatim source excerpt]

                                                                                     q
    For the result, we define the notion of non-trivial reference distributions PN     : a reference distribution
                                                       1
pq : {xi } → △NY is in P N if it assigns mass at least N to the most likely class in the  dataset, i.e., at least 1
data point would be discretized to the most likely class arg maxy i q(y, xi ).13
                                                                   P
\end{ReviewVerbatim}


\paragraph{Expanded PaperInterface specification}
\begin{ReviewVerbatim}
∀ {Sample : Type u_1} {N K : ℕ} (posterior : Sample → Fin N → Fin K → ℝ) (aggregateArgmax : Sample → Fin K)
  (prefAt : Sample → Fin K → ℝ),
  DSWG24DiscretizationBias.ProofBridge.sourcePNq posterior aggregateArgmax prefAt ↔
    DSWG24DiscretizationBias.ProofBridge.referenceSimplexAt prefAt ∧
      (∀ (sample : Sample),
          DSWG24DiscretizationBias.ProofBridge.isAggregateFirstArgmax (posterior sample) (aggregateArgmax sample)) ∧
        ∀ (sample : Sample), 1 / ↑N ≤ prefAt sample (aggregateArgmax sample)
\end{ReviewVerbatim}

\paragraph{Lean proof endpoint (built separately)}\begin{ReviewVerbatim}
DSWG24DiscretizationBias.PaperInterface.nontrivial_reference_family_definition
\end{ReviewVerbatim}

\paragraph{Recorded source-to-Spec screening}
\textbf{Verdict:} \texttt{not recorded}
\\
\textbf{Reason:} not recorded
\reviewmatch{Matches source input}
\noindent\textbf{Reviewer annotation}\par
\reviewerbox
\clearpage
\hypertarget{source-claim-d69492a6cc9555ec}{}
\section*{Review row 11}
\textbf{Source locator:} {\footnotesize\raggedright cited publication, lines 597-\allowbreak{}607\par}
\paragraph{Verbatim source input}
\begin{ReviewVerbatim}
3.1.1    Continuous classifier and discretization rules
Suppose we have a machine learning classifier q (e.g., trained on some other dataset) that makes continuous
probability predictions for each class, given the data. Let q(y, x) denote the probability placed on class y



                                                       10

[form-feed in source extraction]
given data x. For example, the Bayes optimal classifier is
                                               q ∗ (y, x) ≜ Pr(Y = y|X = x).
We will refer to p as the prior and q as our learned posterior. Let the simplex of possible probability vectors
over the support Y be △Y , and let q(x) ∈ △Y be the vector of probabilities (q(y, x))y .
\end{ReviewVerbatim}


\paragraph{Expanded PaperInterface specification}
\begin{ReviewVerbatim}
∀ {X : Type u_1} {Y : Type u_2} [inst : Fintype Y] (q : X → Y → ℝ),
  DSWG24DiscretizationBias.ProofBridge.posteriorSimplex q ↔
    (∀ (x : X), ∑ y, q x y = 1) ∧ (∀ (x : X) (y : Y), 0 ≤ q x y) ∧ ∀ (x : X) (y : Y), q x y ≤ 1
\end{ReviewVerbatim}

\paragraph{Lean proof endpoint (built separately)}\begin{ReviewVerbatim}
DSWG24DiscretizationBias.PaperInterface.posterior_simplex_definition
\end{ReviewVerbatim}

\paragraph{Recorded source-to-Spec screening}
\textbf{Verdict:} \texttt{not recorded}
\\
\textbf{Reason:} not recorded
\reviewmatch{Matches source input}
\noindent\textbf{Reviewer annotation}\par
\reviewerbox
\clearpage
\hypertarget{source-claim-54a3ee9bbb8f086b}{}
\section*{Review row 12}
\textbf{Source locator:} {\footnotesize\raggedright cited publication, lines 635-\allowbreak{}658\par}
\paragraph{Verbatim source input}
\begin{ReviewVerbatim}
3.2     Theoretical Results: Argmax bias and continuous classifier
We first study the relationship between argmax bias and the properties of the continuous classifier. We show
that discretization bias can occur even with unbiased (calibrated) continuous classifiers q, as long as it is
uncertain (has non-zero error).
    The following result holds for all calibrated classifiers. A given classifier q is calibrated when its contin-
uous predictions are correct on average, Pr(Y = y|q(y, x) = c) = c.10 Note that the Bayes optimal classifier
is calibrated.
Theorem 1. Argmax bias depends on predictive uncertainty, i.e., how much information features x provide
about the true class label. Consider calibrated classifier q and the argmax decision rule Dargmax . Let N > K
and consider a reference distribution of either the aggregate posterior or the prior. Then,
  (i). Suppose x provides no information and q(y, x) = Pr(y) for all x, y. Then, amplification bias is
       maximized and all data points are classified as a plurality class:
                                                  (
                                                    1 − Pr(y), if y = arg maxw p(w),
                            BIAS(y, Dargmax ) =
                                                    − Pr(y),      otherwise.
   9 Note that, though related, mean average error of the continuous predictor q(y, x) is not equivalent to expected zero-one

loss of the discrete decisions. The latter, our measure of decision accuracy, also depends on how the continuous predictor is
discretized.                                n                                           o
  10 For all c of non-zero measure: ∀c ∈
                                                 R
                                             c : (x,y) I[q(y, x) = c]f (x, y)d(y, x) > 0 . Further note that calibration implies
Bayesian consistency, i.e., the continuous model cumulatively assigns the proper mass to each class y: EX [q(y, x)] = Pr(y).
\end{ReviewVerbatim}


\paragraph{Expanded PaperInterface specification}
\begin{ReviewVerbatim}
∀ {X : Type u_1} {Y : Type u_2} [inst : MeasurableSpace (X × Y)] [inst_1 : DecidableEq Y]
  (μ : MeasureTheory.Measure (X × Y)) (q : X → Y → ℝ),
  DSWG24DiscretizationBias.ProofBridge.calibrated μ q ↔
    ∀ (y : Y) (s : Set ℝ),
      MeasurableSet s →
        ∫ (xy : X × Y), if q xy.1 y ∈ s then if xy.2 = y then 1 else 0 else 0 ∂μ =
          ∫ (xy : X × Y), if q xy.1 y ∈ s then q xy.1 y else 0 ∂μ
\end{ReviewVerbatim}

\paragraph{Lean proof endpoint (built separately)}\begin{ReviewVerbatim}
DSWG24DiscretizationBias.PaperInterface.calibration_definition
\end{ReviewVerbatim}

\paragraph{Recorded source-to-Spec screening}
\textbf{Verdict:} \texttt{not recorded}
\\
\textbf{Reason:} not recorded
\reviewmatch{Matches source input}
\noindent\textbf{Reviewer annotation}\par
\reviewerbox
\clearpage
\hypertarget{source-claim-861b407dcf52f582}{}
\section*{Review row 13}
\textbf{Source locator:} {\footnotesize\raggedright cited publication, lines 642-\allowbreak{}650\par}
\paragraph{Verbatim source input}
\begin{ReviewVerbatim}
Theorem 1. Argmax bias depends on predictive uncertainty, i.e., how much information features x provide
about the true class label. Consider calibrated classifier q and the argmax decision rule Dargmax . Let N > K
and consider a reference distribution of either the aggregate posterior or the prior. Then,
  (i). Suppose x provides no information and q(y, x) = Pr(y) for all x, y. Then, amplification bias is
       maximized and all data points are classified as a plurality class:
                                                  (
                                                    1 − Pr(y), if y = arg maxw p(w),
                            BIAS(y, Dargmax ) =
                                                    − Pr(y),      otherwise.
\end{ReviewVerbatim}


\paragraph{Expanded PaperInterface specification}
\begin{ReviewVerbatim}
∀ {X : Type u_1} {Y : Type u_2} [inst : Fintype X] [inst_1 : DecidableEq X] [Nonempty X] [inst_3 : Fintype Y]
  [inst_4 : DecidableEq Y] (μ : PMF (X × Y)) (q : X → Y → ℝ) (z y : Y),
  (∀ (x : X) (a : Y), q x a = DSWG24DiscretizationBias.ProofBridge.prior μ a) →
    (∀ (a : Y), DSWG24DiscretizationBias.ProofBridge.prior μ a ≤ DSWG24DiscretizationBias.ProofBridge.prior μ z) →
      (DSWG24DiscretizationBias.Finite.paperBias μ (fun x => z) (DSWG24DiscretizationBias.ProofBridge.prior μ) y =
          if z = y then 1 - DSWG24DiscretizationBias.ProofBridge.prior μ y
          else -DSWG24DiscretizationBias.ProofBridge.prior μ y) ∧
        DSWG24DiscretizationBias.Finite.paperBias μ (fun x => z)
            (DSWG24DiscretizationBias.Finite.paperAggregatePosterior μ q) y =
          if z = y then 1 - DSWG24DiscretizationBias.ProofBridge.prior μ y
          else -DSWG24DiscretizationBias.ProofBridge.prior μ y
\end{ReviewVerbatim}

\paragraph{Lean proof endpoint (built separately)}\begin{ReviewVerbatim}
DSWG24DiscretizationBias.PaperInterface.theorem1i_no_information_bias
\end{ReviewVerbatim}

\paragraph{Recorded source-to-Spec screening}
\textbf{Verdict:} \texttt{not recorded}
\\
\textbf{Reason:} not recorded
\reviewmatch{Matches source input}
\noindent\textbf{Reviewer annotation}\par
\reviewerbox
\clearpage
\hypertarget{source-claim-d8ad786b7cc58022}{}
\section*{Review row 14}
\textbf{Source locator:} {\footnotesize\raggedright cited publication, lines 667-\allowbreak{}673\par}
\paragraph{Verbatim source input}
\begin{ReviewVerbatim}
(iii). More generally, in the intermediate regime, argmax bias is upper bounded by the predictive error of
        the classifier q:

                                                                                                   BIAS(y, Dargmax ) ≤ MAE(q).                                                         (4)

       The bound is tight: there exist FXY and q such that Equation (4) holds with equality for the plurality
       class.
\end{ReviewVerbatim}


\paragraph{Expanded PaperInterface specification}
\begin{ReviewVerbatim}
∀ {X : Type u_1} {Y : Type u_2} [inst : MeasurableSpace X] [inst_1 : MeasurableSpace Y] [MeasurableSingletonClass Y]
  [inst_3 : Fintype Y] [inst_4 : DecidableEq Y] (μ : MeasureTheory.Measure (X × Y)) [MeasureTheory.IsFiniteMeasure μ]
  (q : X → Y → ℝ) (argmaxRule : X → Y) (y : Y),
  Measurable argmaxRule →
    DSWG24DiscretizationBias.ProofBridge.isArgmaxRule q argmaxRule →
      DSWG24DiscretizationBias.ProofBridge.posteriorSimplex q →
        (∀ (y : Y), Measurable fun xy => q xy.1 y) →
          DSWG24DiscretizationBias.ProofBridge.calibrated μ q →
            DSWG24DiscretizationBias.ProofBridge.continuousJointPriorBias μ argmaxRule y ≤
              DSWG24DiscretizationBias.ProofBridge.continuousJointClassifierMAE μ q
\end{ReviewVerbatim}

\paragraph{Lean proof endpoint (built separately)}\begin{ReviewVerbatim}
DSWG24DiscretizationBias.PaperInterface.theorem1iii_argmax_bias_le_mae
\end{ReviewVerbatim}

\paragraph{Recorded source-to-Spec screening}
\textbf{Verdict:} \texttt{not recorded}
\\
\textbf{Reason:} not recorded
\reviewmatch{Matches source input}
\noindent\textbf{Reviewer annotation}\par
\reviewerbox
\clearpage
\hypertarget{source-claim-d6b9878258465151}{}
\section*{Review row 15}
\textbf{Source locator:} {\footnotesize\raggedright cited publication, lines 667-\allowbreak{}673\par}
\paragraph{Verbatim source input}
\begin{ReviewVerbatim}
(iii). More generally, in the intermediate regime, argmax bias is upper bounded by the predictive error of
        the classifier q:

                                                                                                   BIAS(y, Dargmax ) ≤ MAE(q).                                                         (4)

       The bound is tight: there exist FXY and q such that Equation (4) holds with equality for the plurality
       class.
\end{ReviewVerbatim}


\paragraph{Expanded PaperInterface specification}
\begin{ReviewVerbatim}
DSWG24DiscretizationBias.Finite.paperAggregateBias DSWG24DiscretizationBias.Finite.paperTheorem1TightMu
    DSWG24DiscretizationBias.Finite.paperTheorem1TightQ DSWG24DiscretizationBias.Finite.paperTheorem1TightArgmax 0 =
  DSWG24DiscretizationBias.ProofBridge.classifierMAE DSWG24DiscretizationBias.Finite.paperTheorem1TightMu
    DSWG24DiscretizationBias.Finite.paperTheorem1TightQ
\end{ReviewVerbatim}

\paragraph{Lean proof endpoint (built separately)}\begin{ReviewVerbatim}
DSWG24DiscretizationBias.PaperInterface.theorem1iii_tight_binary_example
\end{ReviewVerbatim}

\paragraph{Recorded source-to-Spec screening}
\textbf{Verdict:} \texttt{not recorded}
\\
\textbf{Reason:} not recorded
\reviewmatch{Matches source input}
\noindent\textbf{Reviewer annotation}\par
\reviewerbox
\clearpage
\hypertarget{source-claim-82ddd45159ec7ca3}{}
\section*{Review row 16}
\textbf{Source locator:} {\footnotesize\raggedright cited publication, lines 756-\allowbreak{}758\par}
\paragraph{Verbatim source input}
\begin{ReviewVerbatim}
Theorem 2. Consider Bayes optimal q, and N > K. For all F :
   (i). For every γ, pref there exists a joint decision-making rule that maximizes Oγ (D, q, pref ) in Equa-
        tion (3).
\end{ReviewVerbatim}


\paragraph{Expanded PaperInterface specification}
\begin{ReviewVerbatim}
∀ {ω : Type u_1} {σ : Type u_2} {N K : ℕ} [inst : NeZero K],
  2 ≤ K →
    K < N →
      ∀ (expect : (ω → ℝ) → ℝ),
        EconCSLib.Decision.FiniteLinearExpectation expect →
          ∀ (observedDataset : ω → σ) (trueLabels : ω → Fin N → Fin K) (γ : ℝ) (posterior : σ → Fin N → Fin K → ℝ)
            (fidelityTerm : σ → (Fin N → Fin K) → ℝ),
            (∀ (i : Fin N) (choose : σ → Fin K),
                (expect fun x => if choose (observedDataset x) = trueLabels x i then 1 else 0) =
                  expect fun x => posterior (observedDataset x) i (choose (observedDataset x))) →
              ∃ optRule,
                ∀ (rule : σ → Fin N → Fin K),
                  DSWG24DiscretizationBias.ProofBridge.expectedObjective expect observedDataset trueLabels γ
                      fidelityTerm rule ≤
                    DSWG24DiscretizationBias.ProofBridge.expectedObjective expect observedDataset trueLabels γ
                      fidelityTerm optRule
\end{ReviewVerbatim}

\paragraph{Lean proof endpoint (built separately)}\begin{ReviewVerbatim}
DSWG24DiscretizationBias.PaperInterface.theorem2i_joint_rule_exists
\end{ReviewVerbatim}

\paragraph{Recorded source-to-Spec screening}
\textbf{Verdict:} \texttt{not recorded}
\\
\textbf{Reason:} not recorded
\reviewmatch{Matches source input}
\noindent\textbf{Reviewer annotation}\par
\reviewerbox
\clearpage
\hypertarget{source-claim-cc8ba41c41752422}{}
\section*{Review row 17}
\textbf{Source locator:} {\footnotesize\raggedright cited publication, lines 760\par}
\paragraph{Verbatim source input}
\begin{ReviewVerbatim}
(ii). The argmax decision rule Dargmax maximizes O1 (D, q, pref ), i.e., is accuracy maximizing.
\end{ReviewVerbatim}


\paragraph{Expanded PaperInterface specification}
\begin{ReviewVerbatim}
∀ {ω : Type u_1} {σ : Type u_2} {N K : ℕ} [NeZero K],
  2 ≤ K →
    K < N →
      ∀ (expect : (ω → ℝ) → ℝ),
        EconCSLib.Decision.FiniteLinearExpectation expect →
          ∀ (observedDataset : ω → σ) (trueLabels : ω → Fin N → Fin K) (posterior : σ → Fin N → Fin K → ℝ)
            {decisionRule argmaxRule : σ → Fin N → Fin K},
            (∀ (xs : σ), EconCSLib.Decision.IsPointwiseMax (posterior xs) (argmaxRule xs)) →
              (∀ (i : Fin N) (choose : σ → Fin K),
                  (expect fun x => if choose (observedDataset x) = trueLabels x i then 1 else 0) =
                    expect fun x => posterior (observedDataset x) i (choose (observedDataset x))) →
                EconCSLib.Decision.expectedDecisionAccuracy expect observedDataset trueLabels decisionRule ≤
                  EconCSLib.Decision.expectedDecisionAccuracy expect observedDataset trueLabels argmaxRule
\end{ReviewVerbatim}

\paragraph{Lean proof endpoint (built separately)}\begin{ReviewVerbatim}
DSWG24DiscretizationBias.PaperInterface.theorem2ii_argmax_accuracy_maximizing
\end{ReviewVerbatim}

\paragraph{Recorded source-to-Spec screening}
\textbf{Verdict:} \texttt{not recorded}
\\
\textbf{Reason:} not recorded
\reviewmatch{Matches source input}
\noindent\textbf{Reviewer annotation}\par
\reviewerbox

\end{document}
